% ID: SEL-D0019-OB-E03
% ESTADO: cotejada. Segunda pasada visual completa y borrador compilado revisados.
% Fuente: Selectividad/matematicas_aplicadas/reserva a/Examen-Matemáticas Aplicadas a las Ciencias Sociales.pdf
% SHA-256: 6fed0384ff6fcbaaa7098ab3f35ea5f881d1280502f050ada3a7f13af528392e
% Página: 2; opción B; ejercicio original 3.

Sean $A$, $B$, $C$, $D$, $E$ y $F$ sucesos de un experimento aleatorio.
\begin{enumerate}[label=\alph*)]
\item \textbf{(0.5 puntos)} Se sabe que $P(A)=0.5$, $P(A\cup B)=0.7$ y $P(A\cap B)=0.4$. Halle la probabilidad de que ocurra $B$.
\item \textbf{(1 punto)} Se sabe que $P(C)=0.4$, $P(D)=0.3$ y $P(C\cup D)=0.5$. Halle la probabilidad de que ocurra $C$ sabiendo que no ocurre $D$.
\item \textbf{(1 punto)} Se sabe que los sucesos $E$ y $F$ son independientes, que $P(E)=0.6$ y que $P(F)=0.8$. Calcule la probabilidad de que no ocurra ninguno de los dos sucesos.
\end{enumerate}
